\documentclass{beamer}

\usepackage{minted}
\usetheme{codecentric}

\renewcommand*{\insertdate}{January 26, 2017}

\author{Markus Hauck}
\institute{codecentric AG}

\title{Warm fluffy Interpreters}

\usemintedstyle{trac}

\begin{document}
{
  \usebackgroundtemplate{\includegraphics[width=\paperwidth,height=\paperheight]{background.png}}
  \begin{frame}[plain,noframenumbering]
    \titlepage{}
  \end{frame}
}

\begin{frame}
  \tableofcontents
\end{frame}

\section{Introduction}

\begin{frame}
  \frametitle{Who Am I}
  \begin{itemize}
  \item two years ago, I was sitting in COPL just like you
  \item finished my master at TU Darmstadt October 2015
  \item started to work at codecentric as IT Consultant
  \end{itemize}
\end{frame}

\begin{frame}
  \frametitle{What codecentric does}
  \begin{itemize}
  \item codecentric is about being experts and sharing knowledge (\textbf{both}
    is expected)
  \item we have many people being experts in different topics
  \item Java (performance), Microservices w/o docker, Big Data, Scala, NoSQL,
    \ldots
  \item working model: 4 + 1
  \end{itemize}
\end{frame}

\begin{frame}
  \frametitle{What I Do}
  \begin{itemize}
  \item in university, I already was a fan of expressive programming
    languages
  \item mostly Haskell and Scala, due to love of FP
  \item goal: don't spend the rest of my life writing Java
  \end{itemize}
\end{frame}

\begin{frame}
  \frametitle{What I Do}
  \begin{center}
    \includegraphics[width=0.7\textwidth]{pics/success.jpg}
  \end{center}
\end{frame}

\begin{frame}
  \frametitle{What I Do}
  \begin{itemize}
  \item working for a customer in the online retail market
  \item programming the backend for the webshop
  \item doing Scala Trainings (Lightbend)
  \item speaking at Scala conferences
  \item from time to time: other trainings (e.g. git)
  \end{itemize}
  \vfill{}
  \includegraphics[height=3mm]{pics/codecentric-logo.png}
  \hfill{}
  \includegraphics[height=3mm]{pics/lightbend}
\end{frame}

\begin{frame}
  \frametitle{Warm, Fluffy Interpreters}
  \begin{itemize}
  \item so what for today?
  \item we will see yet another variation on interpreters using\ldots
  \end{itemize}
  \vspace{5mm}
  \begin{minipage}{0.75\linewidth}
    \begin{quote}
      Simon Peyton Jones once commented that the biggest mistake Haskell
      made was to call them ``monads'' instead of ``warm, fluffy things''.
    \end{quote}
  \end{minipage}
  \begin{minipage}{25mm}
    \includegraphics[width=2cm]{pics/spj-snow.jpg}
  \end{minipage}
  \vspace{5mm}
  \begin{itemize}
  \item but first: typeclasses
  \end{itemize}
\end{frame}

\section{Typeclasses}

\begin{frame}
  \frametitle{One Last Thing}
  \begin{itemize}
  \item typeclasses are not in Scala's standard library
  \item by now several good FP libraries that provide typeclasses and
    instances (Scalaz, cats, \ldots)
  \item we will use ``cats'' \url{https://github.com/typelevel/cats}
  \end{itemize}
\end{frame}

\begin{frame}
  \frametitle{Typeclasses}
  \includegraphics[height=10mm]{pics/haskell.png}
  \vspace{3mm}
  \begin{itemize}
  \item design pattern coming from Haskell
  \item another form of ad-hoc polymorphism
  \item can be encoded in Scala using implicits
  \item implicits: values / methods marker such that the compiler will
    try to fill them in if required at call site (compile phase)
  \end{itemize}
  \begin{center}
    \hfill{}
    \includegraphics[height=10mm]{pics/smooth-spiral.png}
  \end{center}
\end{frame}

\begin{frame}[fragile]
  \frametitle{Typeclasses in Scala}
  \begin{itemize}
  \item trait with at least one type parameter defining the available
    methods
  \item implicit value/def on supported types that define the
    behaviour
  \item common but actually optional: syntactic sugar via implicit class
  \end{itemize}
\end{frame}

\begin{frame}[fragile]
  \frametitle{The Monoid Typeclass: Trait}
\begin{minted}[style=trac]{scala}
trait Monoid[A] {
  def empty: A
  def combine(lhs: A, rhs: A): A
}
\end{minted}
\end{frame}

\begin{frame}[fragile]
  \frametitle{The Monoid Typeclass: Instance}
\begin{minted}[style=trac]{scala}
object Monoid {
  def apply[A:Monoid]: Monoid[A] = implicitly

  implicit val plusMonoid: Monoid[Int] = new Monoid {
    def empty: Int = 0
    def combine(lhs: Int, rhs: Int): Int = 0
  }
}
\end{minted}
\end{frame}

\begin{frame}[fragile]
  \frametitle{The Monoid Typeclass: Sugar}
\begin{minted}[style=trac]{scala}
implicit class MonoidOps[A](self: A)(
  implicit M: Monoid[A]
) {
  def |+|(that: A): A = M.combine(self,that)
}
\end{minted}
\end{frame}

\begin{frame}[fragile]
  \frametitle{The Monoid Typeclass: Usage}
\begin{minted}[style=trac]{scala}
scala> import Monoid._
scala> 1 |+| 2
res0: Int = 3
scala> Monoid[Int].empty
res1: Int = 0
scala> Monoid[Int].empty |+| 42
res2: Int = 42
\end{minted}
\end{frame}

\section{Monads}

\begin{frame}
  \frametitle{Monads}
  \begin{itemize}
  \item known in category theory
  \item discovered in Haskell to handle effects
  \item by now very common in FP and also found in Java, Scala, F\#
    and others
  \item for something to ``be a monad'' means to have a (lawful)
    instance of the \texttt{Monad} typeclass
  \end{itemize}
\end{frame}

\begin{frame}[fragile]
  \frametitle{Monad Typeclass}
\begin{minted}[style=trac]{scala}
trait Monad[F[_]] {
  def pure[A](x: A): F[A]
  def flatMap[A, B](fa: F[A])(f: A => F[B]): F[B]
}
\end{minted}
  \begin{itemize}
  \item \texttt{pure}: wrap up a value in a context (boring)
  \item \texttt{flatMap}: all about \textit{control flow}:
    \begin{itemize}
    \item get the \texttt{A} and pass it to \texttt{f}
    \item how to get \texttt{A} is the instance's business
    \end{itemize}
  \end{itemize}
\end{frame}

\begin{frame}[fragile]
  \frametitle{Monad Instance For Option}
\begin{minted}[style=trac]{scala}
implicit val optionMonad: Monad[Option] = new Monad {
  override def pure[A](x: A): Option[A] = Some(x)
  override def flatMap[A,B](fa: Option[A])(
    f: A => Option[B]): Option[B] = {
      fa match {
        case None => None
        case Some(x) => f(x)
      }
  }
}
\end{minted}
\end{frame}

\section{Free Monads \& Interpreters}

\begin{frame}
  \frametitle{Free Monads}
  \begin{itemize}
  \item<1-> ``Free X'' for X in \{Monad, Functor, Monoid, \dots\}
  \item<1-> Idea: strip away any semantics, only keep the AST
  \item<1-> Quiz: what nodes do we have for monads? (what operations?)
  \item<2> pure \& flatMap
  \end{itemize}
\end{frame}

\begin{frame}[fragile]
  \frametitle{Free Monads}
\begin{minted}[style=trac]{scala}
trait Monad[F[_]] {
  def pure[A](x: A): F[A]
  def flatMap[A, B](fa: F[A])(f: A => F[B]): F[B]
}

sealed trait Free[F[_], A]

final case class Pure[F[_], A](x: A)
  extends Free[F, A]

final case class FlatMap[F[_], A](k: F[Free[F, A]])
  extends Free[F, A]
\end{minted}
\end{frame}

\begin{frame}
  \frametitle{Essence of Free Monads}
  \begin{itemize}
  \item known from the current course: build AST, interpret it
  \item up until now: very clumsy syntax
  \item free monads is just another variation of a familiar theme
  \item \textbf{build} the AST via normal map/flatMap/pure operations
  \item \textbf{interpret} the resulting program via interpreter
  \end{itemize}
\end{frame}

\begin{frame}
  \frametitle{Applying Free Monads}
  \begin{center}
    \includegraphics[height=5cm]{pics/free-monad-all.jpg}
  \end{center}
\end{frame}

\begin{frame}[fragile]
  \frametitle{Applying Free Monads}
\begin{minted}[style=trac]{scala}
sealed trait SimpleAlgebra[A]

case class Plus(x: Int, y: Int)
  extends SimpleAlgebra[Int]
case class Mult(x: Int, y: Int)
  extends SimpleAlgebra[Int]

def plus(x: Int, y: Int): Free[SimpleAlgebra, Int] =
  Free.liftF(Plus(x,y))

def mult(x: Int, y: Int): Free[SimpleAlgebra, Int] =
  Free.liftF(Mult(x,y))
\end{minted}
\end{frame}

\begin{frame}
  \frametitle{A More Realistic Scenario}
  \begin{itemize}
  \item Szenario: Our program interacts with Database
  \item<1> for testing we want to stub it
  \item<1> existing solutions: embedded version, in-memory etc.
  \item<2> in some cases still undesirable:
    \begin{itemize}
    \item embedded version requires substantial setup time
    \item makes parallel running tests very resource intensive
    \end{itemize}
  \item<2> better: base your operations on free monads, use different
    interpreters for production and test
  \end{itemize}
\end{frame}

\begin{frame}[fragile]
  \frametitle{Our Program}
\begin{minted}[style=trac]{scala}
def addItem(iid: ItemId, cid: CartId) = for {
  item <- loadItem(iid)
  cart <- loadCart(cid)
  newCart = cart.add(item)
  r <- saveCart(newCart)
} yield r
\end{minted}
\end{frame}

\section{Conclusion}

\begin{frame}
  \frametitle{Conclusion}
\end{frame}

\end{document}
